% =============================================================================
% Annexe A : Extraits de Code
% =============================================================================
\chapter{Extraits de Code}
\label{annexe:code}

\section{Configuration Playwright}

\begin{lstlisting}[language=Java, caption={playwright.config.ts --- Configuration principale}, label={lst:playwright-config}]
import { defineConfig, devices } from '@playwright/test';
import dotenv from 'dotenv';
import path from 'path';

dotenv.config({ path: path.resolve(__dirname, '.env') });

export default defineConfig({
  testDir: '.',
  fullyParallel: true,
  forbidOnly: !!process.env.CI,
  retries: process.env.CI ? 2 : 0,
  workers: process.env.CI ? 2 : undefined,
  timeout: 60_000,

  reporter: [
    ['html', { open: 'never' }],
    ['allure-playwright'],
    ['list'],
  ],

  use: {
    trace: 'on-first-retry',
    screenshot: 'only-on-failure',
    video: 'on-first-retry',
  },

  projects: [
    {
      name: 'admin-auth-setup',
      testDir: './fixtures',
      testMatch: 'admin-auth.setup.ts',
    },
    {
      name: 'admin-panel',
      testDir: './e2e/admin-panel',
      use: {
        ...devices['Desktop Chrome'],
        baseURL: process.env.ADMIN_BASE_URL,
        storageState: 'fixtures/.auth/admin.json',
      },
      dependencies: ['admin-auth-setup'],
    },
    {
      name: 'api-tests',
      testDir: './api',
      use: {
        baseURL: process.env.API_GATEWAY_URL,
      },
    },
  ],
});
\end{lstlisting}

\section{Client API réutilisable}

\begin{lstlisting}[language=Java, caption={api-client.ts --- Client HTTP authentifié}, label={lst:api-client}]
import { APIRequestContext, request } from '@playwright/test';

export class ApiClient {
  private baseUrl: string;
  private token: string | null = null;

  constructor(baseUrl?: string) {
    this.baseUrl = baseUrl
      || process.env.API_GATEWAY_URL
      || 'http://localhost:8080';
  }

  async login(email: string, password: string): Promise<string> {
    const ctx = await request.newContext({
      baseURL: this.baseUrl
    });
    const response = await ctx.post('/api/auth/login', {
      data: { email, password },
    });
    const body = await response.json();
    this.token = body.accessToken || body.token;
    await ctx.dispose();
    return this.token!;
  }

  async get(path: string) {
    const ctx = await this.authenticatedContext();
    const res = await ctx.get(path);
    await ctx.dispose();
    return res;
  }

  async post(path: string, data: unknown) {
    const ctx = await this.authenticatedContext();
    const res = await ctx.post(path, { data });
    await ctx.dispose();
    return res;
  }
}
\end{lstlisting}

\section{Test unitaire backend --- AuthServiceImpl}

\begin{lstlisting}[language=Java, caption={AuthServiceImplTest.java --- Test de connexion}, label={lst:auth-test}]
@ExtendWith(MockitoExtension.class)
class AuthServiceImplTest {

    @Mock private UserRepository userRepository;
    @Mock private JwtTokenProvider jwtTokenProvider;
    @Mock private PasswordEncoder passwordEncoder;
    @InjectMocks private AuthServiceImpl authService;

    @Nested
    @DisplayName("login()")
    class LoginTests {

        @Test
        @DisplayName("Should return AuthResponse for active user")
        void login_activeUser_returnsAuthResponse() {
            User user = new User();
            user.setEmail("test@example.com");
            user.setPassword("hashed");
            user.setStatus(UserStatus.ACTIVE);
            user.setIsEmailVerified(true);

            when(userRepository.findByEmail("test@example.com"))
                .thenReturn(Optional.of(user));
            when(passwordEncoder.matches("raw", "hashed"))
                .thenReturn(true);
            when(jwtTokenProvider.generateToken(user))
                .thenReturn("jwt-token");

            Object result = authService.login(
                "test@example.com", "raw");

            assertThat(result)
                .isInstanceOf(AuthResponse.class);
            verify(jwtTokenProvider).generateToken(user);
        }
    }
}
\end{lstlisting}

\section{Test E2E --- Page Object Model}

\begin{lstlisting}[language=Java, caption={AdminLoginPage --- Pattern Page Object}, label={lst:page-object}]
import { Page, Locator, expect } from '@playwright/test';

export class AdminLoginPage {
  readonly page: Page;
  readonly emailInput: Locator;
  readonly passwordInput: Locator;
  readonly loginButton: Locator;
  readonly errorMessage: Locator;

  constructor(page: Page) {
    this.page = page;
    this.emailInput = page.getByLabel(/email/i);
    this.passwordInput = page.getByLabel(/password/i);
    this.loginButton = page.getByRole('button',
      { name: /sign in|login/i });
    this.errorMessage = page.locator(
      '[class*="error"], mat-error');
  }

  async goto() {
    await this.page.goto('/auth/login');
  }

  async login(email: string, password: string) {
    await this.emailInput.fill(email);
    await this.passwordInput.fill(password);
    await this.loginButton.click();
  }

  async expectDashboardRedirect() {
    await this.page.waitForURL('**/dashboard');
    await expect(this.page).toHaveURL(/dashboard/);
  }
}
\end{lstlisting}

\section{Test de performance k6}

\begin{lstlisting}[language=Java, caption={k6-load-test.js --- Test de charge}, label={lst:k6-load}]
import http from 'k6/http';
import { check, sleep } from 'k6';
import { Rate, Trend } from 'k6/metrics';

const loginDuration = new Trend('login_duration');
const errorRate = new Rate('errors');

export const options = {
  stages: [
    { duration: '30s', target: 10 },
    { duration: '1m',  target: 25 },
    { duration: '2m',  target: 25 },
    { duration: '1m',  target: 50 },
    { duration: '1m',  target: 50 },
    { duration: '30s', target: 0 },
  ],
  thresholds: {
    http_req_duration: ['p(95)<3000'],
    http_req_failed: ['rate<0.05'],
    login_duration: ['p(95)<2000'],
  },
};

export default function () {
  const res = http.post(BASE_URL + '/api/auth/login',
    JSON.stringify({
      email: 'admin@speedline-test.com',
      password: 'TestAdmin123!'
    }),
    { headers: {
      'Content-Type': 'application/json' } }
  );
  loginDuration.add(res.timings.duration);
  check(res, {
    'login 200': (r) => r.status === 200,
  });
  sleep(1);
}
\end{lstlisting}

\section{Test de sécurité --- Injection SQL}

\begin{lstlisting}[language=Java, caption={security.api.spec.ts --- Test d'injection SQL}, label={lst:security}]
test('should handle SQL injection in search', async () => {
  const response = await adminApi.get(
    "/api/orders?search=' OR '1'='1"
  );
  expect([200, 400]).toContain(response.status());
  if (response.status() === 200) {
    const body = await response.json();
    const items = body.content || body;
    expect(Array.isArray(items)).toBeTruthy();
  }
});

test('should handle SQL injection in POST', async ({request}) => {
  const loginRes = await request.post(
    API_BASE + '/api/auth/login', {
    data: {
      email: "admin' OR '1'='1",
      password: "' OR '1'='1",
    },
  });
  expect(loginRes.status()).toBeGreaterThanOrEqual(400);
});
\end{lstlisting}
